% STATUS: written 2026-11-05 (Bennis and Fahid 2017, Boudine and Zerra 2024 read in full)
% STATUS: extended 2026-11-06 (Barajas and Duarte 2026 read in full)
% Sources read in full: Bennis and Fahid 2017 (pp. 1--22, sections 2 to 4);
% Boudine and Zerra 2024 (pp. 459--472, sections 2 and 3); Barajas and Duarte 2026 (pp. 1--2, 6, 13, 15).
% Origins: the trivial extension A x M and its centre, Bennis and Fahid 2017 (p. 4);
% the form of a derivation (Lemma 3.1, p. 8) and the innerness criterion (Theorem 3.12, p. 14);
% the (d1, s1, s2) form over a ring and the prime-ring theorems, Boudine and Zerra 2024
% (Theorem 2.1 p. 461, Theorem 2.3 p. 463, Theorem 3.5 p. 467, Proposition 3.6 p. 468).
\section{Trivial extensions}\label{sec:path-incidence:trivial-extension}

The triangular algebra of Chapter~\ref{ch:matrix-triangular} is a trivial extension
\cite[p.~5]{BennisFahid2017}. A trivial extension puts a ring and a bimodule together
with one rule for the product \cite[p.~4]{BennisFahid2017}. The triangular algebra
is the special case where the module has the form $M=eSf$
\cite[p.~5]{BennisFahid2017}. This section gives the derivations of the
construction and the test for innerness.

Let $A$ be a ring and $M$ an $A$-bimodule. The product of two pairs keeps the first
coordinate and adds a crossed term to the second \cite[p.~4]{BennisFahid2017}.

\begin{definition}[{\cite[p.~4]{BennisFahid2017}}]\label{def:trivial:extension}
Let $A$ be a unital $R$-algebra and $M$ a unital $A$-bimodule. The direct product
$A\times M$, with pairwise addition and the product
\begin{equation*}
(a,m)(b,n)=(ab,\ an+mb)\qquad(a,b\in A,\ m,n\in M),
\end{equation*}
is a unital algebra, called the \emph{trivial extension} of $A$ by $M$ and written
$A\ltimes M$.
\end{definition}

Every linear map on the extension splits into four parts
\cite[p.~8]{BennisFahid2017}. Two parts act on $A$ and $M$ separately, and two carry
$M$ to $A$ and $A$ to $M$ \cite[p.~8]{BennisFahid2017}.

\begin{lemma}[{\cite[Lemma~3.1, p.~8]{BennisFahid2017}}]\label{lem:trivial:form}
Every linear map $D\colon A\ltimes M\to A\ltimes M$ has the form
\begin{equation*}
D(a,m)=(D_{A}(a)+T(m),\ D_{M}(a)+S(m))\qquad(a\in A,\ m\in M)
\end{equation*}
for some linear maps $D_{A}\colon A\to A$, $D_{M}\colon A\to M$, $T\colon M\to A$ and
$S\colon M\to M$. We write $D=(D_{A}+T,D_{M}+S)$.
Moreover, $D$ is a derivation if and only if $D_{A}$ and $D_{M}$ are derivations,
$T\in E(M)$ and $S$ is a module generalized $D_{A}$-derivation.
\end{lemma}

The notion of $E(M)$ in Lemma~\ref{lem:trivial:form} needs its own definition
\cite[p.~7]{BennisFahid2017}. It is the set of bimodule maps whose sum over two
arguments vanishes \cite[p.~7]{BennisFahid2017}. A module generalized derivation
satisfies the product rule on both sides, with the same derivation $d$ of $A$ on each
side \cite[p.~7]{BennisFahid2017}.

\begin{definition}[{\cite[p.~7]{BennisFahid2017}}]\label{def:trivial:egeneral}
Let $E(M)$ be the $R$-submodule of $\mathrm{Hom}_{A-A}(M,A)$ consisting of all
bimodule homomorphisms $f$ such that $f(m)n+mf(n)=0$ for all $m,n\in M$. A linear
map $S\colon M\to M$ is a \emph{module generalized $d$-derivation} for some
$d\in\mathrm{Der}(A)$ if
\begin{equation*}
S(am)=aS(m)+d(a)m,\qquad S(ma)=S(m)a+md(a)
\end{equation*}
for all $m\in M$ and $a\in A$.
\end{definition}

Innerness of the whole map is decided by its four parts
\cite[p.~14]{BennisFahid2017}. One part must vanish, and the two parts on the module
must be inner (Theorem~\ref{thm:trivial:inner}) \cite[p.~14]{BennisFahid2017}.

\begin{theorem}[{\cite[Theorem~3.12, p.~14]{BennisFahid2017}}]\label{thm:trivial:inner}
Consider a derivation $D\colon A\ltimes M\to A\ltimes M$ of the form
$D=(D_{A}+T,D_{M}+S)$, where $D_{A}$, $T$, $D_{M}$ and $S$ are as in
Lemma~\ref{lem:trivial:form}. Then the following assertions hold.
\begin{enumerate}[label=(\arabic*)]
\item If $D$ is inner, then $T=0$ and both $S$ and $D_{M}$ are inner.
\item The converse holds if $\mathrm{l.Ann}_{A}(M)\cap\mathrm{r.Ann}_{A}(M)=0$.
\end{enumerate}
\end{theorem}

When the module is the ring itself, the criterion takes an exact form
\cite[p.~14]{BennisFahid2017}. The two module parts then agree \cite[p.~14]{BennisFahid2017}.

\begin{corollary}[{\cite[Corollary~3.13, p.~14]{BennisFahid2017}}]\label{cor:trivial:self}
Assume $M=A$ and consider a derivation $D\colon A\ltimes M\to A\ltimes M$ of the
form $D=(D_{A}+T,D_{M}+S)$ with the parts as in Lemma~\ref{lem:trivial:form}. Then
$D$ is inner if and only if $T=0$, $D_{M}$ is inner, $S=D_{A}$ and $S$ is inner.
Precisely, the inner derivations on $A\ltimes A$ are only of the form
$(d_{1},d_{0}+d_{1})$, where $d_{0}$ and $d_{1}$ are inner derivations on $A$.
\end{corollary}

Boudine and Zerra give the same description over a general ring with two
projections \cite[p.~461]{BoudineZerra2024}. Write $d=(d_{1},d_{2})$ for the two
components, and write $s_{1}(m)=d_{2}((0,m))$ and $s_{2}(r)=d_{2}((r,0))$
\cite[p.~461]{BoudineZerra2024}. Their conditions are the three product rules of
Bennis and Fahid, written out.

\begin{theorem}[{\cite[Theorem~2.1, p.~461]{BoudineZerra2024}}]\label{thm:trivial:components}
With the notation above, the map $d$ is a derivation on $R\propto M$ if and only if
the following statements hold:
\begin{enumerate}[label=(\arabic*)]
\item $d_{1}$ is a $(\pi_{1},\pi_{1})$-derivation;
\item $s_{2}$ is a derivation;
\item $s_{1}(rm)=r\,s_{1}(m)+d_{1}(r,0)m$ and $s_{1}(mr)=s_{1}(m)r+m\,d_{1}(r,0)$
for every $r\in R$ and $m\in M$.
\end{enumerate}
\end{theorem}

Their innerness test lists four conditions, one for each part
\cite[p.~463]{BoudineZerra2024}. The list matches Corollary~\ref{cor:trivial:self} in
the case $M=R$, and both are the case $R=M$ of
Theorem~\ref{thm:trivial:inner-ring}.

\begin{theorem}[{\cite[Theorem~2.3, p.~463]{BoudineZerra2024}}]\label{thm:trivial:inner-ring}
With the notation above, the derivation $d$ is the inner derivation on
$R\propto M$ induced by $(r,m)\in R\propto M$ if and only if the following
statements hold:
\begin{enumerate}[label=(\arabic*)]
\item $d_{1}\circ i_{2}=0$;
\item $d_{1}\circ i_{1}$ is the inner derivation on $R$ induced by $r$;
\item $s_{1}$ is the inner derivation on $M$ induced by $r$;
\item $s_{2}$ is the inner derivation on $R$ into $M$ induced by $m$.
\end{enumerate}
\end{theorem}

Over a prime ring the classical theorems of Herstein, Posner and Bell carry over to
the extension \cite[p.~467]{BoudineZerra2024}. The first one kills the part of the
derivation that stays on $R$ (Theorem~\ref{thm:trivial:herstein}) \cite[p.~467]{BoudineZerra2024}.

\begin{theorem}[{\cite[Theorem~3.5, p.~467]{BoudineZerra2024}}]\label{thm:trivial:herstein}
Let $R$ be a noncommutative prime ring with $\mathrm{char}(R)\neq 2$, $M$ an
$R$-bimodule, and $d$ a derivation on $R\propto M$. If
$d((r,m))d((r',m'))=d((r',m'))d((r,m))$ for every $(r,m),(r',m')\in R\propto M$,
then $d_{1}\circ i_{1}=0$ and $s_{1}\in\mathrm{End}_{R}(M)$. Furthermore, for every
$m,m'\in M$, $[s_{2}(m),d_{1}((0,m'))]=[d_{1}((0,m)),s_{2}(m')]$.
\end{theorem}

The centralizing form gives the corresponding statement for the theorem of Posner (Proposition~\ref{prop:trivial:posner})
\cite[p.~468]{BoudineZerra2024}.\begin{proposition}[{\cite[Proposition~3.6, p.~468]{BoudineZerra2024}}]\label{prop:trivial:posner}
Let $R$ be a noncommutative prime ring, $M$ an $R$-bimodule such that
$Z_{M}(R)=M$, and $d$ a derivation on $R\propto M$. Then $d$ is centralizing if and
only if, for every $(r,m)\in R\propto M$, the following statements hold:
\begin{enumerate}[label=(\arabic*)]
\item $d_{1}\circ i_{1}=0$;
\item $d_{1}((0,m))r-rd_{1}((0,m))\in Z_{R}(R)$.
\end{enumerate}
\end{proposition}

By Theorem~\ref{thm:trivial:components}, the extension keeps the derivations of its
first factor and adds a crossed part. That pattern returns once more with a graph in
place of the module, and Section~\ref{sec:path-incidence:leavitt} gives the graph
case.

The construction repeats for a finite family of modules, and the derivations
become a sequence of maps \cite[p.~2]{BarajasDuarte2026}. Anderson, Bennis,
Fahid and Shaiea defined the $n$-trivial extension of a ring by modules
$M_1,\dots,M_n$ \cite[p.~2]{BarajasDuarte2026}. Barajas and Duarte ask for the
maps that stand to it as derivations stand to the one-module case
\cite[p.~2]{BarajasDuarte2026}. When all the modules are the same algebra, the
answer is the Hasse--Schmidt sequence of
Section~\ref{sec:higher-multi:hasse-schmidt}, which is the higher derivation
under its other name \cite[p.~2]{BarajasDuarte2026}.
For different algebras they define a multi Hasse--Schmidt derivation, a sequence
of additive maps whose Leibniz rule carries the maps $\varphi_{ij}$ between the
modules \cite[Definition~3.1, p.~6]{BarajasDuarte2026}. The $n$-trivial
extension is then a quotient of a polynomial ring, exactly as the one-module
extension is the ring of dual numbers \cite[Theorem~4.2, p.~13]{BarajasDuarte2026}.
The multi Hasse--Schmidt derivations with first term the identity correspond to
the algebra maps into that quotient \cite[Corollary~4.3, p.~15]{BarajasDuarte2026}.

For a Poisson algebra $A$ and a left Poisson module $M$ a Poisson structure is put on the trivial extension $A \oplus M$, and its Poisson and Hamiltonian derivations are determined together with the first Poisson cohomology group~\cite{WlPoissonDerivationsAnd2019}.

Every nonlinear Jordan triple derivable map on a 2-torsion free $*$-type trivial extension algebra is described in standard form~\cite{WlNonlinearJordanTriple2024}.

Every higher Jordan triple derivation on a 2-torsion free $*$-type trivial extension algebra is described in standard form~\cite{WlHigherJordanTriple2024}.

Fei, Wang, Lu and Zhang show that every higher Jordan triple derivation of a 2-torsion free $*$-type trivial extension algebra is a sum of a higher derivation and a higher anti-derivation, and apply this to characterise the higher Jordan triple derivations of triangular algebras~\cite{WlHigherJordanTripleDerivationsO2024}.
